% ============================================================
%  Chapter 5: Discussion
% ============================================================

\chapter{Discussion}
\label{chap:discussion}

% ============================================================
%  Section 5.1
% ============================================================
\section{Contextualising the TFPnP–FBPConvNet comparison}
\label{sec:why-fbpconvnet}

The reproduction confirms~\citet{Wei2022TFPnP}'s central claims. The learned policy improves on hand-calibrated hyperparameters by $+2.96$\,dB at 5\% noise (\Cref{sec:headline}), and TFPnP degrades gracefully across the noise range (\Cref{sec:cross-noise}). The one direct disagreement with the paper is the FBPConvNet–TFPnP ranking at 5\% noise, which reverss, but can be explained by two factors.

The first is an in-domain training asymmetry. In~\citet{Wei2022TFPnP}, TFPnP's denoiser is trained on BSD400 and its policy on PASCAL VOC, both natural-image datasets. FBPconv's training data is not stated explicitly, but the paper's relativedly flat FBPconv PSNRs across noise levels ($23.00, 22.94, 22.86$\,dB) are consistent with out-of-domain training. In this reproduction, FBPConvNet is trained end-to-end on LIDC, while TFPnP's policy is also LIDC-trained but the denoiser remains the natural-image DRUNet (\Cref{tab:implementation-differences}). FBPConvNet is fully in-domain, yet TFPnP is only partially so, and at low noise the fully in-domain method extracts fine anatomical detail from the FBP reconstruction more effectively than the natural-image denoiser can.

The second is architectural. TFPnP overtakes FBPConvNet at 10\% noise on all three metrics (\Cref{fig:metrics_vs_noise}), with a PSNR drop across the range ($-1.17$\,dB) roughly a third of FBPConvNet's ($-3.85$\,dB). Part of this reflects a training-distribution effect, as FBPConvNet was trained at fixed 5\% noise while TFPnP saw the mixed regime, but the crossover on SSIM and HaarPSI is unlikely to be explained by that alone. Repeated denoising passes appear to remove sinogram noise progressively in a way a single forward pass cannot. The paper's out-of-domain comparison obscured both effects, flattering TFPnP relative to FBPconv on the averaged ranking.

% ============================================================
%  Section 5.2
% ============================================================
\section{Does an in-domain denoiser close the gap?}
\label{sec:in-domain-gap}

If FBPConvNet's advantage at low noise comes from its in-domain training, a natural extension is to give TFPnP the same advantage by training the denoiser on CT rather than natural images. The CT-DRUNet extension experiment (\Cref{sec:ct-drunet}) yields two findings and one negative result.

The first finding is that in the Fixed PnP-ADMM setting, the CT-trained DRUNet outperforms the natural-image DRUNet at every noise level, with the advantage growing with noise (\Cref{tab:fixed_pnp_ct_vs_natural}). This result stands independently of what happens with the full TFPnP pipeline.

The second is that the CT-DRUNet operates at a fundamentally different $\sigma$ scale than the natural-image DRUNet. The natural DRUNet peaks at $\sigma = 1.5$ (\Cref{fig:sigma_sweep_natural}), the CT-DRUNet at $\sigma = 8.0$ (\Cref{fig:sigma_sweep_ct}). Domain-specific training therefore does not just refine the denoiser's ability to reject noise, but shifts the operating point at which the denoiser is most useful, with direct consequences for how policy $\sigma$ ranges must be configured when swapping denoisers within TFPnP.

The negative result is that retraining the TFPnP policy against the CT-DRUNet (\texttt{run\_11}) with a widened $\sigma$ action range failed. The critic loss diverged and the policy $\sigma$ output collapsed shortly afterwards. The obvious explanation is the action range itself, as widening $\sigma$ from $[1, 5]$ to $[1, 15]$ enlarges the space of reachable ADMM states, and therefore the range of values the critic must learn to estimate. A second explanation is our own value network, as Table~1 of~\citet{Wei2022TFPnP} specifies that the critic uses the policy's feature extractor with batch normalisation replaced by weight normalisation and ReLU replaced by TReLU. Our critic removes batch normalisation, as the paper does, but puts nothing in its place, using unnormalised convolutions and standard ReLU activations (\Cref{sec:tfpnp-policy}). Weight normalisation is introduced specifically to stabilise training, but our critic has neither and diverges.

Taken together, these results suggest that the in-domain denoiser advantage is real and reproducible in the simpler Fixed PnP-ADMM setting, but that we were unable to translate it into the RL-based TFPnP pipeline. Whether that reflects a sensitivity of TFPnP's RL formulation to action-range choices, or our own substitution of the value network, remains open.
% ============================================================
%  Section 5.3 — On the choice of evaluation metrics
% ============================================================
\section{On the choice of evaluation metrics}
\label{sec:metrics-discussion}

PSNR is the primary metric throughout this report, and was chosen for direct comparison to~\citet{Wei2022TFPnP}. However, PSNR captures only pixel-wise mean squared error, which does not always correlate with the visual quality in medical imaging. PSNR weights every pixel equally and is insensitive to the structure of errors, so a slightly blurred reconstruction that smooths away texture can score identically to one that preserves texture but adds a small amount of noise, although a radiologist would rate the two very differently. SSIM partially addresses the second issue by comparing local luminance, contrast, and structural patterns rather than only pixel intensities. HaarPSI extends this further through a Haar wavelet decomposition, and has been shown to correlate more strongly with radiological assessment than either PSNR or SSIM~\citep{breger2025study}. 

% ============================================================
%  Section 5.4 — On direct comparison with the paper
% ============================================================
\section{On direct comparison with~\citet{Wei2022TFPnP}}
\label{sec:direct-comparison}

The absolute PSNR values reported in this project cannot be directly compared with those in~\citet{Wei2022TFPnP}, due to the differences detailed below.

First, the PSNR data-range convention differs, with this report following LION's per-image \texttt{data\_range} $= x_{\text{gt}}.\text{max}()$ in native $\mu$ units, whereas the paper does not state its convention but most likely uses \texttt{data\_range} $= 1.0$ on $[0, 1]$-normalised images. This alone accounts for a $\sim 5$\,dB shift in reported PSNR for the same mean squared error (\Cref{sec:geometry-noise}), affecting direct comparisons but not within-experiment rankings.

Second, the paper reports on 50 randomly selected lung slices from the COVID-19 CT lung and infection segmentation dataset, and this project reports on the full 389-image LIDC-IDRI test set. The two datasets have different image statistics, noise characteristics, and anatomy distributions, so absolute PSNR is not directly comparable even at identical noise levels.

Third, the sinogram noise convention differs, as while both the paper and this reproduction follow a percentage-of-signal-standard-deviation convention, the paper's scaling is not documented in detail, so replication is inexact (\Cref{sec:geometry-noise}), but LION's own more physical Poisson-Gaussian noise model was not used to mitigate this issue as much as possible.

Fourth, several implementation choices differ, consolidated in~\Cref{tab:implementation-differences}. The denoiser is trained with MSE loss (\texttt{train\_drunet\_ct.py}) rather than the $L_1$ loss used by~\citet{Wei2022TFPnP} and~\citet{zhang2021plug}, chosen because PSNR is monotone in MSE and so aligns the denoiser with the RL reward $\Delta\text{PSNR}$, at the cost of slightly more oversmoothing. The denoiser batch size was reduced to 8 instead of 32 to fit within the CSD3 GPU memory budget alongside the RL training runs. The policy learning rates are split into three groups ($10^{-4}$ critic, $3\times 10^{-5}$ shared policy features and termination head, $10^{-6}$ parameter head; \texttt{train\_tfpnp.py}) rather than the paper's two-group schedule, reflecting the higher sensitivity of $\pi_2$'s model-based gradient through the differentiable ADMM environment, and no decay schedule was used in this reproduction. Finally, the $\mu$ action range is $[10, 100]$ rather than the paper's $[0.01, 1]$, not as a hyperparameter choice but as a consequence of the $z$-step gradient normalisation (\Cref{sec:admm-gradient}), since LION's adjoint produces unnormalised backprojections $\sim 100\times$ larger than TorchRadon's, so $\mu$ has to be $100\times$ larger to balance them.

Because of these differences, the reproduction's central claim is the within-experiment finding that TFPnP outperforms Fixed PnP-ADMM by $+2.96$\,dB at 5\% noise, not the absolute agreement of numbers between the two studies. The FBPConvNet–TFPnP ranking inversion at 5\% noise is a similarly within-experiment finding, and its interpretation depends on the in-domain vs out-of-domain training asymmetry discussed in~\Cref{sec:why-fbpconvnet} rather than on absolute PSNR values.

Beyond these study-level differences, two limitations of the comparison itself should be noted. First, three of the paper's six baselines (RED-CNN, RPGD, and LPD) are either absent from LION or would require training beyond the project timeline, so our comparison covers only a subset of the paper's landscape. Second, all training runs used a single random seed due to compute constraints, so the reported numbers reflect one RL trajectory rather than an ensemble. Combined with the four factors above, this means the reproduction's central claims are best interpreted as within-experiment findings rather than as absolute number-for-number reproductions of the paper's results.

% ============================================================
%  Section 5.5 — Broader implications
% ============================================================
\section{Broader implications}
\label{sec:implications}

Three broader observations emerge from this reproduction.

First,~\citet{Wei2022TFPnP}'s fair-comparison methodology, in which all learning-based methods are evaluated on CT without domain-matched training data, may obscure the true capabilities of end-to-end methods on specific medical imaging domains. Allowing FBPConvNet to train in-domain, substantially changes its performance relative to TFPnP. This does not invalidate the paper's core contribution, but it does suggest that PnP methods need to be compared against domain-adapted baselines to inform practical deployment.

Second, TFPnP's tuning-free claim is nuanced. It is accurate at inference time, but the training pipeline itself requires careful tuning of $\sigma$/$\mu$ action ranges, learning rates, and reward functions. The reward-shaping ablation (\Cref{sec:reward-shaping}) shows meaningful differences across reward configurations, and the \texttt{run\_11} failure illustrates that swapping in a different denoiser requires extensive re-engineering of the RL training pipeline. Tuning-free at inference does not mean tuning-free during development.

Finally, the same experimental structure, comparing in-domain and out-of-domain denoisers within a plug-and-play framework, could be applied to other PnP algorithms such as HQS, PnP-PGM, and RED. Whether the in-domain denoiser advantage transfers to these methods is an open question that this reproduction's infrastructure would support with some additional work.